\chapter{Répertoire des graphies}\label{ch:rep_graphies}

Les corrections graphiques signalées par les éditeurs (notation ⸢~⸣ dans le catalogue) permettent un relevé des écarts entre la graphie gravée et la norme. Ce relevé dépend des choix des éditeurs, qui ne signalent pas toutes les graphies non classiques de la même manière : il ne mesure pas la fréquence absolue des phénomènes, mais leur visibilité dans l’édition (chapitre~\ref{ch:langue}).

Pour chaque correspondance attestée au moins deux fois (lettre attendue, lettre déclarée gravée), on donne le nombre d’occurrences et de notices, la répartition par période selon l’année de classement, des exemples tirés du texte développé et la liste complète des notices. Les correspondances attestées une seule fois figurent à la fin du tableau~\ref{tab:graphies} et dans l’index des graphies.

{\small
% généré par scripts/04_generated.py
\subsection*{B gravé pour \textit{v} (bêtacisme)}
\noindent Occurrences : 72 dans 52 notices ; par période (ensemble A, année de classement) : 340–359 : 13 ; 360–374 : 20 ; 375–389 : 17 ; 390–403 : 22.\par
\noindent Exemples (forme développée, lettre attendue entre demi-crochets) : \lat{⸢v⸣ixit} (4) ; \lat{⸢v⸣ixit} (6) ; \lat{vi⸢v⸣it} (22) ; \lat{vi⸢v⸣o} (24) ; \lat{Esu⸢v⸣iae} (26) ; \lat{⸢v⸣iginti} (34) ; \lat{⸢v⸣ix⸢i⸣t} (40) ; \lat{⸢v⸣ixit} (47) ; \lat{exi⸢v⸣it} (56) ; \lat{Ne⸢v⸣i⸢t⸣} (63) ; \lat{Ne⸢v⸣i⸢tt⸣a} (72) ; \lat{No⸢v⸣embres} (83).\par
\noindent{\small Notices : 4, 6, 22, 24, 26, 34, 40, 47, 56, 63, 72, 83, 89, 90, 91, 106, 112, 113, 118, 119, 120, 124, 144, 172, 184, 194, 207, 208, 215, 217, 234, 235, 261, 272, 281, 301, 313, 332, 347, 352, 365, 366, 374, 402, 407, 408, 417, 440, 441, 444, 468, 475.}\par\medskip
\subsection*{V gravé pour \textit{o} (confusion o/u)}
\noindent Occurrences : 44 dans 42 notices ; par période (ensemble A, année de classement) : 340–359 : 3 ; 360–374 : 5 ; 375–389 : 10 ; 390–403 : 26.\par
\noindent Exemples (forme développée, lettre attendue entre demi-crochets) : \lat{marit⸢o⸣} (32) ; \lat{ann⸢o⸣s} (54) ; \lat{ann⸢o⸣s} (55) ; \lat{ann⸢o⸣s} (69) ; \lat{ann⸢o⸣s} (94) ; \lat{ann⸢o⸣s} (104) ; \lat{c⸢o⸣mpare} (130) ; \lat{ann⸢o⸣s} (161) ; \lat{ann⸢o⸣s} (191) ; \lat{ann⸢o⸣s} (193) ; \lat{tinian⸢o⸣} (209) ; \lat{ann⸢o⸣s} (225).\par
\noindent{\small Notices : 32, 54, 55, 69, 94, 104, 130, 161, 191, 193, 209, 225, 275, 281, 288, 289, 299, 320, 322, 323, 327, 336, 346, 347, 353, 359, 365, 374, 375, 376, 379, 400, 419, 421, 427, 431, 443, 466, 470, 473, 475, 478.}\par\medskip
\subsection*{I gravé pour \textit{e} (confusion e/i)}
\noindent Occurrences : 43 dans 35 notices ; par période (ensemble A, année de classement) : 340–359 : 6 ; 360–374 : 7 ; 375–389 : 7 ; 390–403 : 23.\par
\noindent Exemples (forme développée, lettre attendue entre demi-crochets) : \lat{dec⸢e⸣ssit} (4) ; \lat{d⸢e⸣possitus} (25) ; \lat{mens⸢e⸣s} (28) ; \lat{mens⸢e⸣s} (34) ; \lat{mens⸢e⸣s} (55) ; \lat{ctobr⸢e⸣s} (60) ; \lat{Septembr⸢e⸣s} (80) ; \lat{mens⸢e⸣s} (84) ; \lat{mens⸢e⸣s} (94) ; \lat{mens⸢e⸣s} (124) ; \lat{⸢e⸣t} (151) ; \lat{menss⸢e⸣s} (157).\par
\noindent{\small Notices : 4, 25, 28, 34, 55, 60, 80, 84, 94, 124, 151, 157, 191, 193, 260, 261, 287, 306, 337, 347, 353, 359, 365, 366, 367, 381, 389, 390, 395, 411, 423, 431, 441, 443, 446.}\par\medskip
\subsection*{I gravé pour \textit{t} (i pour t (lecture))}
\noindent Occurrences : 18 dans 15 notices ; par période (ensemble A, année de classement) : 340–359 : 1 ; 360–374 : 5 ; 375–389 : 8 ; 390–403 : 4.\par
\noindent Exemples (forme développée, lettre attendue entre demi-crochets) : \lat{e⸢t⸣} (37) ; \lat{Ne⸢v⸣i⸢t⸣} (63) ; \lat{Criscen⸢t⸣inae} (132) ; \lat{len⸢t⸣e} (136) ; \lat{e⸢t⸣} (153) ; \lat{e⸢t⸣} (164) ; \lat{vixi⸢t⸣} (186) ; \lat{e⸢t⸣} (231) ; \lat{vixi⸢t⸣} (245) ; \lat{ce⸢t⸣ius} (296) ; \lat{eri⸢t⸣} (306) ; \lat{vixi⸢t⸣} (308).\par
\noindent{\small Notices : 37, 63, 132, 136, 153, 164, 186, 231, 245, 296, 306, 308, 333, 371, 433.}\par\medskip
\subsection*{E gravé pour \textit{i} (confusion e/i)}
\noindent Occurrences : 17 dans 16 notices ; par période (ensemble A, année de classement) : 340–359 : 4 ; 360–374 : 5 ; 375–389 : 1 ; 390–403 : 7.\par
\noindent Exemples (forme développée, lettre attendue entre demi-crochets) : \lat{iscrips⸢i⸣t} (27) ; \lat{⸢v⸣ix⸢i⸣t} (40) ; \lat{al⸢i⸣} (50) ; \lat{s⸢i⸣t} (54) ; \lat{quiesc⸢i⸣t} (70) ; \lat{⸢i⸣n} (76) ; \lat{merent⸢i⸣} (100) ; \lat{⸢q⸣u⸢i⸣} (111) ; \lat{incompara⸢b⸣il⸢i⸣s} (125) ; \lat{ostend⸢i⸣s} (306) ; \lat{fec⸢i⸣t} (353) ; \lat{fec⸢i⸣t} (364).\par
\noindent{\small Notices : 27, 40, 50, 54, 70, 76, 100, 111, 125, 306, 353, 364, 368, 388, 391, 443.}\par\medskip
\subsection*{I gravé pour \textit{y} (graphie de y)}
\noindent Occurrences : 13 dans 13 notices ; par période (ensemble A, année de classement) : 340–359 : 1 ; 360–374 : 2 ; 375–389 : 4 ; 390–403 : 6.\par
\noindent Exemples (forme développée, lettre attendue entre demi-crochets) : \lat{A⸢c⸣ind⸢y⸣no} (1) ; \lat{neof⸢y⸣to} (145) ; \lat{neof⸢y⸣ta} (172) ; \lat{of⸢y⸣to} (233) ; \lat{S⸢y⸣agrio} (238) ; \lat{neof⸢y⸣tus} (276) ; \lat{f⸢y⸣tus} (308) ; \lat{Eut⸢y⸣chius} (352) ; \lat{neof⸢y⸣tus} (391) ; \lat{Eut⸢y⸣chiano} (417) ; \lat{Eut⸢y⸣chiano} (418) ; \lat{t⸢y⸣chiano} (419).\par
\noindent{\small Notices : 1, 145, 172, 233, 238, 276, 308, 352, 391, 417, 418, 419, 424.}\par\medskip
\subsection*{C gravé pour \textit{g} (c pour g)}
\noindent Occurrences : 13 dans 13 notices ; par période (ensemble A, année de classement) : 340–359 : 2 ; 360–374 : 1 ; 375–389 : 5 ; 390–403 : 5.\par
\noindent Exemples (forme développée, lettre attendue entre demi-crochets) : \lat{sarcofa⸢g⸣o} (19) ; \lat{Ni⸢g⸣riniano} (35) ; \lat{vir⸢g⸣ini} (99) ; \lat{Au⸢g⸣usto} (201) ; \lat{Sya⸢g⸣rio} (221) ; \lat{⸢g⸣loria} (231) ; \lat{coniu⸢g⸣i} (304) ; \lat{ve⸢g⸣etans} (306) ; \lat{beni⸢g⸣nus} (352) ; \lat{Eu⸢g⸣enio} (353) ; \lat{Au⸢g⸣usto} (354) ; \lat{Eu⸢g⸣e} (359).\par
\noindent{\small Notices : 19, 35, 99, 201, 221, 231, 304, 306, 352, 353, 354, 359, 366.}\par\medskip
\subsection*{N gravé pour \textit{m} (n pour m)}
\noindent Occurrences : 12 dans 11 notices ; par période (ensemble A, année de classement) : 340–359 : 3 ; 360–374 : 3 ; 375–389 : 3 ; 390–403 : 3.\par
\noindent Exemples (forme développée, lettre attendue entre demi-crochets) : \lat{Dece⸢m⸣bres} (23) ; \lat{co⸢m⸣} (36) ; \lat{co⸢m⸣pari} (46) ; \lat{No⸢v⸣e⸢m⸣bres} (119) ; \lat{Dece⸢m⸣bres} (126) ; \lat{co⸢m⸣pare} (184) ; \lat{co⸢m⸣parasse} (217) ; \lat{co⸢m⸣parem!} (275) ; \lat{co⸢m⸣pare} (365) ; \lat{Septe⸢m⸣bres} (419) ; \lat{co⸢m⸣paratus} (424).\par
\noindent{\small Notices : 23, 36, 46, 119, 126, 184, 217, 275, 365, 419, 424.}\par\medskip
\subsection*{V gravé pour \textit{b} (bêtacisme)}
\noindent Occurrences : 9 dans 9 notices ; par période (ensemble A, année de classement) : 340–359 : 0 ; 360–374 : 3 ; 375–389 : 1 ; 390–403 : 5.\par
\noindent Exemples (forme développée, lettre attendue entre demi-crochets) : \lat{super⸢b⸣ia} (85) ; \lat{o⸢b⸣iit} (123) ; \lat{incompara⸢b⸣il⸢i⸣s} (125) ; \lat{⸢b⸣ene} (276) ; \lat{al⸢b⸣a} (340) ; \lat{s⸢b⸣i} (357) ; \lat{Pro⸢b⸣ino} (380) ; \lat{acer⸢b⸣us} (382) ; \lat{acer⸢b⸣o} (470).\par
\noindent{\small Notices : 85, 123, 125, 276, 340, 357, 380, 382, 470.}\par\medskip
\subsection*{B gravé pour \textit{V}}
\noindent Occurrences : 7 dans 5 notices ; par période (ensemble A, année de classement) : 340–359 : 1 ; 360–374 : 5 ; 375–389 : 0 ; 390–403 : 1.\par
\noindent Exemples (forme développée, lettre attendue entre demi-crochets) : \lat{⸢V⸣eneris} (55) ; \lat{⸢V⸣arroniano} (90) ; \lat{⸢V⸣alentiniano} (97) ; \lat{⸢V⸣alentiniano} (128) ; \lat{⸢V⸣eneris} (416).\par
\noindent{\small Notices : 55, 90, 97, 128, 416.}\par\medskip
\subsection*{S gravé pour \textit{x} (s pour x)}
\noindent Occurrences : 6 dans 6 notices ; par période (ensemble A, année de classement) : 340–359 : 1 ; 360–374 : 2 ; 375–389 : 3 ; 390–403 : 0.\par
\noindent Exemples (forme développée, lettre attendue entre demi-crochets) : \lat{vi⸢x⸣it} (13) ; \lat{vi⸢x⸣it} (130) ; \lat{vi⸢x⸣it} (132) ; \lat{vi⸢x⸣it} (272) ; \lat{vi⸢x⸣it} (302) ; \lat{vi⸢x⸣it} (304).\par
\noindent{\small Notices : 13, 130, 132, 272, 302, 304.}\par\medskip
\subsection*{I gravé pour \textit{l}}
\noindent Occurrences : 6 dans 5 notices ; par période (ensemble A, année de classement) : 340–359 : 0 ; 360–374 : 0 ; 375–389 : 3 ; 390–403 : 3.\par
\noindent Exemples (forme développée, lettre attendue entre demi-crochets) : \lat{Vita⸢l⸣ini} (186) ; \lat{C⸢l⸣earcho} (266) ; \lat{So⸢l⸣is} (287) ; \lat{A⸢l⸣exandra} (333) ; \lat{consu⸢l⸣ibus} (359).\par
\noindent{\small Notices : 186, 266, 287, 333, 359.}\par\medskip
\subsection*{E gravé pour \textit{a} (e pour a)}
\noindent Occurrences : 5 dans 4 notices ; par période (ensemble A, année de classement) : 340–359 : 1 ; 360–374 : 0 ; 375–389 : 3 ; 390–403 : 1.\par
\noindent Exemples (forme développée, lettre attendue entre demi-crochets) : \lat{Salli⸢a⸣} (28) ; \lat{qu⸢a⸣m} (214) ; \lat{I⸢a⸣nuarias} (279) ; \lat{I⸢a⸣nua} (466).\par
\noindent{\small Notices : 28, 214, 279, 466.}\par\medskip
\subsection*{V gravé pour \textit{y} (graphie de y)}
\noindent Occurrences : 5 dans 5 notices ; par période (ensemble A, année de classement) : 340–359 : 0 ; 360–374 : 0 ; 375–389 : 2 ; 390–403 : 3.\par
\noindent Exemples (forme développée, lettre attendue entre demi-crochets) : \lat{S⸢y⸣a} (223) ; \lat{S⸢y⸣agrio} (239) ; \lat{S⸢y⸣mmacho} (338) ; \lat{S⸢y⸣mmacho} (339) ; \lat{S⸢y⸣mmaco} (340).\par
\noindent{\small Notices : 223, 239, 338, 339, 340.}\par\medskip
\subsection*{O gravé pour \textit{q}}
\noindent Occurrences : 4 dans 4 notices ; par période (ensemble A, année de classement) : 340–359 : 0 ; 360–374 : 1 ; 375–389 : 1 ; 390–403 : 2.\par
\noindent Exemples (forme développée, lettre attendue entre demi-crochets) : \lat{⸢q⸣u⸢i⸣} (111) ; \lat{⸢q⸣ui} (231) ; \lat{⸢q⸣ui} (388) ; \lat{⸢q⸣ui} (418).\par
\noindent{\small Notices : 111, 231, 388, 418.}\par\medskip
\subsection*{VM gravé pour \textit{o}}
\noindent Occurrences : 4 dans 2 notices ; par période (ensemble A, année de classement) : 340–359 : 0 ; 360–374 : 0 ; 375–389 : 1 ; 390–403 : 3.\par
\noindent Exemples (forme développée, lettre attendue entre demi-crochets) : \lat{marit⸢o⸣} (291) ; \lat{fili⸢o⸣} (402).\par
\noindent{\small Notices : 291, 402.}\par\medskip
\subsection*{Q gravé pour \textit{c} (q pour c)}
\noindent Occurrences : 3 dans 3 notices ; par période (ensemble A, année de classement) : 340–359 : 1 ; 360–374 : 0 ; 375–389 : 2 ; 390–403 : 0.\par
\noindent Exemples (forme développée, lettre attendue entre demi-crochets) : \lat{A⸢c⸣ind⸢y⸣no} (1) ; \lat{⸢q⸣uices⸢c⸣uet} (184) ; \lat{⸢i⸣a⸢c⸣et} (287).\par
\noindent{\small Notices : 1, 184, 287.}\par\medskip
\subsection*{C gravé pour \textit{K}}
\noindent Occurrences : 3 dans 3 notices ; par période (ensemble A, année de classement) : 340–359 : 1 ; 360–374 : 1 ; 375–389 : 1 ; 390–403 : 0.\par
\noindent Exemples (forme développée, lettre attendue entre demi-crochets) : \lat{⸢K⸣alendas} (9) ; \lat{⸢K⸣alendas} (128) ; \lat{⸢K⸣alendas} (279).\par
\noindent{\small Notices : 9, 128, 279.}\par\medskip
\subsection*{QV gravé pour \textit{c}}
\noindent Occurrences : 3 dans 3 notices ; par période (ensemble A, année de classement) : 340–359 : 0 ; 360–374 : 1 ; 375–389 : 0 ; 390–403 : 2.\par
\noindent Exemples (forme développée, lettre attendue entre demi-crochets) : \lat{quies⸢c⸣et} (59) ; \lat{⸢c⸣it} (404) ; \lat{⸢c⸣enti} (405).\par
\noindent{\small Notices : 59, 404, 405.}\par\medskip
\subsection*{D gravé pour \textit{t}}
\noindent Occurrences : 3 dans 3 notices ; par période (ensemble A, année de classement) : 340–359 : 0 ; 360–374 : 0 ; 375–389 : 1 ; 390–403 : 2.\par
\noindent Exemples (forme développée, lettre attendue entre demi-crochets) : \lat{reliqui⸢t⸣} (182) ; \lat{feci⸢t⸣} (319) ; \lat{A⸢t⸣tic⸢o⸣} (402).\par
\noindent{\small Notices : 182, 319, 402.}\par\medskip
\subsection*{O gravé pour \textit{a}}
\noindent Occurrences : 3 dans 3 notices ; par période (ensemble A, année de classement) : 340–359 : 0 ; 360–374 : 0 ; 375–389 : 0 ; 390–403 : 3.\par
\noindent Exemples (forme développée, lettre attendue entre demi-crochets) : \lat{Fravvitt⸢a⸣} (448) ; \lat{Fravitt⸢a⸣} (449) ; \lat{t⸢a⸣} (451).\par
\noindent{\small Notices : 448, 449, 451.}\par\medskip
\subsection*{V gravé pour \textit{i}}
\noindent Occurrences : 2 dans 2 notices ; par période (ensemble A, année de classement) : 340–359 : 1 ; 360–374 : 1 ; 375–389 : 0 ; 390–403 : 0.\par
\noindent Exemples (forme développée, lettre attendue entre demi-crochets) : \lat{Arb⸢i⸣tione} (44) ; \lat{Ar⸢i⸣ntheo} (157).\par
\noindent{\small Notices : 44, 157.}\par\medskip
\subsection*{BE gravé pour \textit{vi}}
\noindent Occurrences : 2 dans 2 notices ; par période (ensemble A, année de classement) : 340–359 : 0 ; 360–374 : 1 ; 375–389 : 0 ; 390–403 : 1.\par
\noindent Exemples (forme développée, lettre attendue entre demi-crochets) : \lat{e⸢vi⸣tta} (66) ; \lat{⸢vi⸣t} (458).\par
\noindent{\small Notices : 66, 458.}\par\medskip
\subsection*{V gravé pour \textit{n}}
\noindent Occurrences : 2 dans 2 notices ; par période (ensemble A, année de classement) : 340–359 : 0 ; 360–374 : 1 ; 375–389 : 1 ; 390–403 : 0.\par
\noindent Exemples (forme développée, lettre attendue entre demi-crochets) : \lat{a⸢n⸣nos} (85) ; \lat{co⸢n⸣} (309).\par
\noindent{\small Notices : 85, 309.}\par\medskip
\subsection*{M gravé pour \textit{n}}
\noindent Occurrences : 2 dans 2 notices ; par période (ensemble A, année de classement) : 340–359 : 0 ; 360–374 : 2 ; 375–389 : 0 ; 390–403 : 0.\par
\noindent Exemples (forme développée, lettre attendue entre demi-crochets) : \lat{i⸢n⸣} (121) ; \lat{i⸢n⸣} (150).\par
\noindent{\small Notices : 121, 150.}\par\medskip
\subsection*{VX gravé pour \textit{XV}}
\noindent Occurrences : 2 dans 2 notices ; par période (ensemble A, année de classement) : 340–359 : 0 ; 360–374 : 1 ; 375–389 : 1 ; 390–403 : 0.\par
\noindent Exemples (forme développée, lettre attendue entre demi-crochets) : \lat{⸢XV⸣} (128) ; \lat{⸢XV⸣} (198).\par
\noindent{\small Notices : 128, 198.}\par\medskip
\subsection*{O gravé pour \textit{u} (confusion o/u)}
\noindent Occurrences : 2 dans 2 notices ; par période (ensemble A, année de classement) : 340–359 : 0 ; 360–374 : 1 ; 375–389 : 1 ; 390–403 : 0.\par
\noindent Exemples (forme développée, lettre attendue entre demi-crochets) : \lat{cons⸢u⸣} (143) ; \lat{Moero⸢b⸣a⸢u⸣dis} (299).\par
\noindent{\small Notices : 143, 299.}\par\medskip
\subsection*{O gravé pour \textit{i}}
\noindent Occurrences : 2 dans 2 notices ; par période (ensemble A, année de classement) : 340–359 : 0 ; 360–374 : 1 ; 375–389 : 1 ; 390–403 : 0.\par
\noindent Exemples (forme développée, lettre attendue entre demi-crochets) : \lat{Modest⸢i⸣} (156) ; \lat{Clearch⸢i⸣} (275).\par
\noindent{\small Notices : 156, 275.}\par\medskip
\subsection*{Y gravé pour \textit{e}}
\noindent Occurrences : 2 dans 2 notices ; par période (ensemble A, année de classement) : 340–359 : 0 ; 360–374 : 0 ; 375–389 : 2 ; 390–403 : 0.\par
\noindent Exemples (forme développée, lettre attendue entre demi-crochets) : \lat{M⸢e⸣robaude} (198) ; \lat{M⸢e⸣robaude} (303).\par
\noindent{\small Notices : 198, 303.}\par\medskip
\subsection*{I gravé pour \textit{o}}
\noindent Occurrences : 2 dans 2 notices ; par période (ensemble A, année de classement) : 340–359 : 0 ; 360–374 : 0 ; 375–389 : 1 ; 390–403 : 1.\par
\noindent Exemples (forme développée, lettre attendue entre demi-crochets) : \lat{Ant⸢o⸣nio} (231) ; \lat{A⸢t⸣tic⸢o⸣} (402).\par
\noindent{\small Notices : 231, 402.}\par\medskip
\subsection*{U gravé pour \textit{o}}
\noindent Occurrences : 2 dans 2 notices ; par période (ensemble A, année de classement) : 340–359 : 0 ; 360–374 : 0 ; 375–389 : 2 ; 390–403 : 0.\par
\noindent Exemples (forme développée, lettre attendue entre demi-crochets) : \lat{ann⸢o⸣s} (260) ; \lat{ann⸢o⸣s} (282).\par
\noindent{\small Notices : 260, 282.}\par\medskip
\subsection*{B gravé pour \textit{r}}
\noindent Occurrences : 2 dans 2 notices ; par période (ensemble A, année de classement) : 340–359 : 0 ; 360–374 : 0 ; 375–389 : 1 ; 390–403 : 1.\par
\noindent Exemples (forme développée, lettre attendue entre demi-crochets) : \lat{p⸢r⸣idie} (308) ; \lat{anno⸢r⸣um} (371).\par
\noindent{\small Notices : 308, 371.}\par\medskip
\subsection*{I gravé pour \textit{r}}
\noindent Occurrences : 2 dans 2 notices ; par période (ensemble A, année de classement) : 340–359 : 0 ; 360–374 : 0 ; 375–389 : 0 ; 390–403 : 2.\par
\noindent Exemples (forme développée, lettre attendue entre demi-crochets) : \lat{p⸢r⸣idie} (446) ; \lat{do⸢r⸣mit} (464).\par
\noindent{\small Notices : 446, 464.}\par\medskip

}
